\section{Introduction}

\subsection{Global Timber Trade, CITES Enforcement, and Timber Forensics}
Illegal logging and the associated trade in illicit timber pose severe threats to global biodiversity, accelerate deforestation, and undermine sustainable forest management worldwide~\cite{dormontt2015}. To curb illegal exploitation, international agreements such as the Convention on International Trade in Endangered Species of Wild Fauna and Flora (CITES) regulate the commercial trade of protected timber species listed under CITES Appendix II~\cite{cites}. Effective enforcement of these regulations requires fast, reliable, and legally admissible species identification at border checkpoints and along supply chains. 

Traditional timber identification heavily relies on laboratory-based forensic techniques, including wood anatomical analysis by trained xylotomists, DNA barcoding, and stable isotope ratio analysis~\cite{dormontt2015,wiedenhoeft2011}. Although laboratory forensic methods offer high diagnostic accuracy, they suffer from inherent scalability bottlenecks: they are time-consuming, expensive, destructive to wood samples, and require highly specialized equipment and taxonomic expertise that are rarely available at field inspection sites~\cite{ravindran2022}. Consequently, there is an urgent demand for automated, portable, non-destructive, and scalable visual identification systems capable of assisting customs inspectors in real-time timber verification.

\subsection{Fine-Grained Visual Recognition in Wood Anatomy}
Computer vision and deep learning have emerged as promising avenues for automated wood species identification, enabling rapid screening based on macroscopic cross-sectional imagery~\cite{woodreview,wu2021}. Xylotomical identification relies on subtle anatomical structures embedded in cross-sections, such as vessel pore arrangements, axial parenchyma distribution, ray widths, and growth ring boundaries~\cite{wiedenhoeft2011}. 

From a computer vision perspective, wood identification constitutes a classic fine-grained visual classification (FGVC) problem~\cite{wah2011,woodreview}. Taxonomically close species within the same genus (e.g., \textit{Afzelia} or \textit{Dalbergia}) frequently exhibit minimal inter-class anatomical divergence. Conversely, individual trees within the same species display substantial intra-class morphological variability due to environmental conditions, tree age, radial sampling position, and sectioning angles. Consequently, models must learn highly nuanced anatomical representation spaces capable of suppressing superficial surface variations while accentuating diagnostic taxonomic traits.

\subsection{Data Generation Structure and Source-Level Leakage}
A critical yet frequently overlooked vulnerability in computer-vision-based wood identification lies in the data generation process. In practice, building wood image datasets involves harvesting multiple physical specimens (blocks or cores) from individual trees. During image acquisition, a single physical specimen is typically photographed repeatedly under varying magnifications, section planes, lighting conditions, or sensor angles.

As a result, an image dataset is rarely a set of independent observations. Multiple images originating from the identical physical specimen inherit shared contextual artifacts: identical wood grain patterns, unique surface saw marks, background discoloration, camera sensor noise, and specific lighting gradients. When such datasets are partitioned using conventional \emph{random image-level splits}, images from the exact same physical specimen end up distributed across both training and test partitions. 

Under random image-level partitioning, a deep neural network can easily achieve near-perfect test scores by memorizing source-specific background shortcuts or surface artifacts rather than learning transferable biological structures~\cite{kaufman2012,kapoor2023}. This systemic flaw, known as \emph{source-level data leakage} or \emph{group leakage}, creates an illusion of high predictive accuracy that catastrophically collapses when the model is deployed on truly unseen specimens in field environments~\cite{roberts2021,varoquaux2022}.

\subsection{Problem Statement and Formalization}
Let $\mathcal{X}=\{x_i\}_{i=1}^{M}$ represent a dataset of $M$ macroscopic wood images and $\mathcal{S}=\{S_j\}_{j=1}^{N}$ represent the set of $N$ underlying physical specimens. Each image $x_i$ is associated with a ground-truth species label $y_i \in \{1, \dots, C\}$ and is linked to its physical source specimen via a mapping function $s(x_i) \in \mathcal{S}$.

To evaluate true generalization to novel physical samples, model assessment must be strictly conducted on physical specimens that were completely withheld during model development. Therefore, dataset partitioning must enforce strict group disjointness at the specimen level:
\begin{equation}
\mathcal{S}_{\mathrm{train}}\cap\mathcal{S}_{\mathrm{val}} =
\mathcal{S}_{\mathrm{train}}\cap\mathcal{S}_{\mathrm{test}} =
\mathcal{S}_{\mathrm{val}}\cap\mathcal{S}_{\mathrm{test}} = \varnothing.
\label{eq:group_disjoint}
\end{equation}
The corresponding image subsets $\mathcal{D}_{\mathrm{train}}, \mathcal{D}_{\mathrm{val}}, \mathcal{D}_{\mathrm{test}}$ are subsequently derived from these specimen partitions:
\begin{equation}
\mathcal{D}_{\mathrm{split}}=\{x_i \in \mathcal{X} : s(x_i) \in \mathcal{S}_{\mathrm{split}}\}, \quad \text{split} \in \{\mathrm{train}, \mathrm{val}, \mathrm{test}\}.
\label{eq:train_from_specimen}
\end{equation}

Condition~(\ref{eq:group_disjoint}) guarantees specimen-level isolation provided that provenance identifiers are complete and accurate. It is essential to emphasize that specimen-level disjointness does not require images in different splits to be distant in feature space; distinct physical specimens of the same species may naturally possess biological similarity. High embedding similarity across splits should be interpreted as an audit signal triggering provenance verification, rather than definitive proof of data leakage.

\subsection{Methodological Gap}
The vast majority of existing studies in automated wood identification concentrate almost exclusively on model architectures, backbone selection, or specialized loss functions~\cite{fabijanska2021,wu2021}. While architectural innovations are vital, they address only one segment of the machine learning pipeline. When specimen metadata, session logs, and source provenance are omitted or unverified during data collection, reproducing data splits or validating sample independence becomes virtually impossible~\cite{kapoor2023}. Furthermore, standard group splits can still yield unbalanced class distributions or subtle representation overlap if not properly audited.

This research addresses this gap by expanding the focus from isolated model optimization to establishing a comprehensive, auditable data pipeline for fine-grained visual recognition on unseen physical specimens. Our proposed methodology encompasses the complete data lifecycle: defining source specimen units, recording standardized provenance metadata, enforcing structured manifest governance, conducting pre-split data auditing, and implementing specimen-disjoint partitioning prior to model training and evaluation.

\subsection{Contributions}
The primary contributions of this paper are fourfold:
\begin{itemize}
    \item We propose the \emph{Specimen-Centric Data Protocol} (SCDP), a general data governance framework for image datasets with repeated physical observations, establishing standardized guidelines for collection, hierarchical storage, manifest tracking, auditing, and leakage-free evaluation.
    \item We develop \emph{Class-wise Embedding-Guided Specimen Split} (CEGS-Split), an automated group-disjoint partitioning algorithm that preserves specimen-level isolation, maintains per-species balance, and utilizes frozen feature embeddings to audit cross-split pairs for subtle source overlap.
    \item We quantify the impact of source-level data leakage on the benchmark S3 wood dataset (18 species across 5 genera), demonstrating a significant drop in k-nearest neighbor (KNN) accuracy when transitioning from random image splits to specimen-disjoint splits ($95.81\%$ vs. $90.87\%$).
    \item We establish a unified benchmark comparing classification baselines, core metric learning loss functions, and self-supervised learning algorithms under strict SCDP governance, highlighting the necessity of auditable data protocols in scientific machine learning.
\end{itemize}

\section{Related Work}

\subsection{Wood Species Identification with Computer Vision}
Early automated wood identification systems relied on manual feature extraction techniques, such as gray-level co-occurrence matrices (GLCM) and local binary patterns (LBP), applied to micro-photographs or macro-scans~\cite{woodreview}. Over the past decade, deep convolutional neural networks (CNNs) and vision transformers have revolutionized the domain by automatically extracting hierarchical feature representations from macroscopic wood surfaces~\cite{fabijanska2021,wu2021}. While these models demonstrate high classification accuracy on controlled benchmarks, their performance often degrades in field settings due to domain shifts in lighting, magnification, and unverified data partitioning protocols~\cite{ravindran2022}.

\subsection{Data Governance and Leakage-Free Evaluation}
Data leakage occurs when information from outside the training set unduly influences model development, producing unrealistically optimistic performance metrics~\cite{kaufman2012}. In domains characterized by repeated measurements—such as medical imaging (multiple scans per patient) or remote sensing—failing to split data at the patient or geographic level is a major driver of the machine learning reproducibility crisis~\cite{roberts2021,varoquaux2022,kapoor2023}. These principles directly transfer to wood species identification: the physical specimen (or tree) is the fundamental unit of independence that must remain isolated throughout evaluation.

\subsection{Representation Learning for Fine-Grained Recognition}
Deep metric learning structures feature embedding spaces such that intra-class samples cluster closely while inter-class samples are pushed apart~\cite{musgrave2020}. Loss functions such as ArcFace, Multi-Similarity, Circle Loss, and Supervised Contrastive Learning (SupCon) have achieved state-of-the-art results in fine-grained retrieval tasks~\cite{arcface,multisimilarity,circleloss,supcon}. Similarly, self-supervised learning (SSL) frameworks—including SimCLR, BYOL, SimSiam, and Barlow Twins—learn robust representations without explicit label supervision~\cite{simclr,byol,simsiam,barlow}. In this study, we evaluate metric learning and SSL as downstream representation tools operating strictly within our auditable data protocol.
